\section{Introdução}

Modelos de linguagem de grande porte (\textit{Large Language Models}, LLMs) são
pré-treinados sobre vastos corpora de texto de propósito geral. Embora isso lhes
confira fluência ampla, o conhecimento adquirido tende a ser raso em domínios
especializados e em variedades regionais da língua, pouco representados nos dados
de origem. Adaptar esses modelos a um domínio específico, com custo computacional
acessível, é um problema central tanto na pesquisa quanto na aplicação prática de
Processamento de Linguagem Natural~\cite{gururangan2020dont}.

Este relatório documenta três experimentos de adaptação do modelo
\texttt{Qwen2.5-0.5B}~\cite{qwen2025} ao domínio público piauiense, conduzidos no
âmbito da disciplina de Tópicos em Inteligência Artificial. Cada experimento
corresponde a uma questão do trabalho e a uma técnica distinta de treinamento:

\begin{itemize}[leftmargin=2.2em]
  \item \textbf{Questão 01 -- Pré-Treino Continuado.} Adaptação não supervisionada
        do modelo ao corpus DOMPI-2025 (Diário Oficial dos Municípios do Piauí),
        avaliando a qualidade da modelagem de linguagem antes e depois do treino.
  \item \textbf{Questão 02 -- Pós-Treino com SFT.} Ajuste fino supervisionado
        (\textit{Supervised Fine-Tuning}) sobre pares de pergunta e resposta
        gerados a partir do dataset docentesDC, atualizando todos os parâmetros do
        modelo.
  \item \textbf{Questão 03 -- Pós-Treino com LoRA/QLoRA.} Repetição do experimento
        anterior por meio de adaptação de baixo posto (\textit{Low-Rank
        Adaptation}), que atualiza menos de 1\% dos parâmetros, comparando custo e
        qualidade frente ao ajuste fino completo.
\end{itemize}

As três técnicas são avaliadas pelas mesmas métricas quantitativas de modelagem de
linguagem: perplexidade, entropia cruzada e acurácia de previsão de
tokens~\cite{jelinek1977perplexity}. A hipótese geral é que o treinamento, em todas
as suas formas, reduz a incerteza do modelo sobre o texto do domínio e melhora sua
capacidade de previsão.

\subsection{Estado dos experimentos e escopo do relatório}

A Questão 01 foi executada de ponta a ponta em uma GPU NVIDIA T4 (Kaggle), com
resultados completos antes e depois do treino. As Questões 02 e 03 tiveram a
geração dos dados, o carregamento dos modelos e a avaliação \textit{baseline} (antes
do treino) concluídos; o laço de treinamento, contudo, foi interrompido em uma
tentativa de execução local por restrição de tempo de hardware. Por essa razão, este
relatório apresenta para as Questões 02 e 03 a metodologia completa e os resultados
\textit{baseline} reais já coletados, deixando as métricas pós-treino sinalizadas
como \emph{em andamento}. Essa decisão preserva a honestidade dos números
reportados: nenhuma métrica é estimada ou fabricada. A
Seção~\ref{sec:consideracoes} detalha o caminho de conclusão dessas duas questões.
